\documentclass[10pt,a4paper]{amsart}
\usepackage[margin=19mm]{geometry}
\usepackage{amsmath,amssymb,mathtools}
\usepackage[scr=rsfs,cal=euler]{mathalfa}
\usepackage{xfrac}
\usepackage[unicode,hidelinks,bookmarks=false]{hyperref}
\newcommand{\ie}{i.e.\ }
\newcommand{\cf}{cf.\ }
\newcommand{\resp}{respectively\ }
\newcommand{\Subsection}{\subsection}
\providecommand{\nobreakdash}{\nobreak\hbox{-}}
\setlength{\emergencystretch}{2em}
\multlinegap0pt
\usepackage{bm}
\usepackage{bbm}
\usepackage{dsfont}
\usepackage{xspace}
\usepackage{cancel}
\usepackage{mathtools}
\usepackage{amsthm}
\makeatletter
\@ifundefined{prop}{\newtheorem{prop}{Proposition}}{}
\@ifundefined{remark}{\newtheorem{remark}{Remark}}{}
\makeatother
\newcommand{\Z}{\mathbb Z}
\newcommand{\grM}{\operatorname{gr}^{\mathscr M}}
\title[Presenting the Zassenhaus Lie algebra by the Magnus Lie alge]{Presenting the Zassenhaus Lie algebra by the Magnus Lie algebra}
\author{Ettore Marmo, David Riley, Thomas Weigel}
\date{Source: arXiv:2511.13379v1 (CC BY 4.0)}
\begin{document}
\maketitle

\noindent\emph{Source and licence.} Presenting the Zassenhaus Lie algebra by the Magnus Lie algebra. Version arXiv:2511.13379v1 (CC BY 4.0), distributed under the Creative Commons Attribution 4.0 International licence, as recorded in the arXiv licence metadata for this exact version and shown on the abstract page https://arxiv.org/abs/2511.13379v1. Full rights notice: licenses/arxiv-2511.13379-CC-BY-4.0.txt.\par\smallskip

\noindent\textit{Editorial scope.} Counted unit: the Prop whose source label is \texttt{prop:basis} in the article (source0.tex lines 434--446, source label \texttt{prop:basis}), delivered together with the article's own complete original proof of it (source0.tex lines 447--475, one \texttt{proof} environment of 29 lines). No mathematics is written, repaired or re-derived: statement and proof are the article's own text line for line. Required context retained uncounted: rem:explicit.filtration (lines 330--338); prop:hall (lines 340--346). Every definition, display and label of those lines is retained unchanged. Omitted from this package and not redistributed: 1--329, 339--339, 347--433, 476--623. Nothing inside a retained excerpt is omitted or altered: every retained body is delivered exactly as the article's lines are written, so no omission receipt of any kind accompanies this package. Citations: the retained excerpts carry no citation command at all, so no bibliography entry of the article is transcribed and no reference is redistributed. Reference adaptation: none was needed and none was made; every reference inside the delivered unit resolves inside it, so no unresolved reference or citation is produced. Printed slips kept literal: no printed word, symbol, hypothesis or number of the counted statement or of its proof was altered. Layout and class: the article's own class is \texttt{amsart} with packages including graphicx, amsmath, amssymb, amsthm, xy, mathrsfs, tikz-cd, xcolor, mathtools, url; the wrapper is the lane's pinned \texttt{amsart} editorial wrapper at 10pt on A4, which restates the theorem-like environments the retained text needs and adds the article's own macro definitions that the retained text needs (\texttt{\textbackslash Z}, \texttt{\textbackslash grM}), with extra wrapper packages (bm, bbm, dsfont, xspace, cancel, mathtools). The delivered numbering is the unit's own. Assets: no retained span contains a figure environment, a graphics-inclusion command, a PDF-page inclusion, a table or a TikZ picture; no image file is copied into this package, the package declares no asset dependency and no image file is redistributed with it. Licence: version arXiv:2511.13379v1 of this article is distributed under the Creative Commons Attribution 4.0 International licence, as recorded in the arXiv licence metadata for this exact version and shown on the abstract page https://arxiv.org/abs/2511.13379v1. A full rights notice is delivered at licenses/arxiv-2511.13379-CC-BY-4.0.txt. Characters: every retained excerpt is pure ASCII, so the delivered engine, XeLaTeX with its default Latin Modern text font, reports no missing character.
\begin{remark}\label{rem:explicit.filtration}
    Let $n, i \geq 1$ and define $j(n, i)$ to be the least positive integer $j$ such that $ip^{j} \geq n$. Then, for all $i \leq k \leq n$, one has:
    \begin{equation}                        D_{n, k}(G) = \prod_{i=k}^n \gamma_i(G)^{p^{j(n, i)}} %\gamma_{k+1}(G)^{p^{j(n, k+1)}} \cdots \gamma_n(G)^{p^{j(n, n)}}
    \end{equation}
    in particular, the following inclusions hold:
    \[
        D_n(G) = D_{n, 1}(G) \geq D_{n, 2}(G) \geq \cdots \geq D_{n, n}(G) = \gamma_n(G).
    \] 
\end{remark}

\begin{prop}\label{prop:hall}
    Let $G$ be a group, $i \geq 1$ and $j \geq 0$, the following hold:
    \begin{enumerate}
        \item $(xy)^{p^j} \equiv x^{p^j} y^{p^j} \mod D_{ip^j, i+1}(G)$ for every $x, y \in \gamma_i(G)$;
        \item $[x^{p^j}, y] \in D_{ip^j, i+1}(G)$ for every $x \in \gamma_i(G)$ and $y \in G$.
    \end{enumerate}
\end{prop}

\begin{prop}\label{prop:basis}
    Let $G$ be a nilpotent $\gamma$-free (pro-$p$) group of nilpotency class $c$. For every $n \geq 1$, one can pick a $\Z$-basis (resp. $\Z_p$-basis) $\{x_{n, \lambda} \gamma_{n+1}(G)\}_{\lambda \in \Lambda_n}$ of the free abelian (pro-$p$) group $\grM_n(G) = \gamma_n(G)/\gamma_{n+1}(G)$ where the indexing set $\Lambda_n$ is totally ordered. Then, the following hold:

    \begin{enumerate}
        \item For every $g \in G$ there exist unique elements $\alpha_{n, \lambda} = \alpha_{n, \lambda}(g)$ in $\Z$ (resp. $\Z_p$) with $n \in \{1, \dots, c\}$ and $\lambda \in \Lambda_n$ such that: 
        \[
            g = \prod_{n = 1}^c \prod_{\lambda \in \Lambda_n} x_{n, \lambda}^{\alpha_{n, \lambda}}
        \]
        where the indices in the product are ordered lexicographically (i.e. $(n, \lambda) \prec (n', \lambda')$ if either $n < n'$ or $n = n'$ and $\lambda < \lambda'$ in $\Lambda_n$).

        \item Whenever $g \in D_{n,k}(G)$ for some $n \geq k$, one has that $p^{j(n, i)}$ divides $\alpha_{i, \lambda}(g)$ for all $k \leq i \leq n$.
    \end{enumerate}
\end{prop}

\begin{proof}
    Since $G$ is nilpotent of class $c$, then the graded Lie algebra $\grM_\bullet(G)$ associated to the lower central series is a \emph{finite} direct sum:
    \[
        \grM_\bullet(G) = G/\gamma_2(G) \oplus \cdots \oplus \gamma_{c-1}(G)/\gamma_c(G) \oplus \gamma_c(G)
    \]
    whose terms are all free abelian since $G$ is $\gamma$-free. So the elements $\{ x_{n, \lambda}\gamma_{n+1}(G) \in \grM_n(G) \mid n = 1, \dots, c; \, \lambda \in \Lambda_n\}$ form a $\Z$-basis (resp. $\Z_p$ basis) for $\grM_n(G)$.
    Hence for every $g \in G$ and $1 \leq n \leq c$ there exist unique coefficients $\alpha_{n, \lambda} = \alpha_{n, \lambda}(g)$ such that
    \[
        g = \prod_{n=1}^c\prod_{n, \lambda} x_{i, \lambda}^{\alpha_{n, \lambda}}
    \]
\noindent
    If $y \in \gamma_n(G)$, for each $1 \leq n \leq j$, we can express the image of both elements $y$ and $y^{p^j}$ in the quotient $\grM_n(G)$ as follows:
    \begin{align*}
        y &\equiv \prod\nolimits_{\lambda \in \Lambda_n} x_{n, \lambda}^{\alpha_{n, \lambda}} \mod \gamma_{n+1}(G), \\
        y^{p^j} &\equiv \prod\nolimits_{\lambda \in \Lambda_n} x_{n, \lambda}^{\beta_{n, \lambda}} \mod \gamma_{n+1}(G)
    \end{align*}
    where the integers $\alpha_{n, \lambda} = \alpha_{n, \lambda}(y)$ and $\beta_{n, \lambda} = \beta_{n, \lambda}(y^{p^j})$ for $\lambda \in \Lambda_n$ are uniquely determined as we proved in the first claim.
    From Proposition \ref{prop:hall} (1) we also find
    \[
        y^{p^j} \equiv \left( \prod\nolimits_{\lambda\in\Lambda_n} x_{n,\lambda}^{\alpha_n,\lambda}\right)^{p^j} \equiv \prod\nolimits_{\lambda\in\Lambda_n} x_{n,\lambda}^{p^j\alpha_n,\lambda} \mod \gamma_{n+1}(G)
    \]
    so, by uniqueness, we must have $\beta_{n, \lambda} = p^j \alpha_{n, \lambda}$ for all $\lambda \in \Lambda_n$, i.e. $p^j$ divides each coefficient $\beta_{n, \lambda}$.
\noindent
    Finally, by Remark \ref{rem:explicit.filtration}, any $g \in D_{n, k}(G)$ can be expressed as a product
    \[
        g = y_k^{p^{j(n, k)}} y_{k-1}^{p^{j(n, k-1)}} \cdots y_{n}^{p^{j(n, n)}}
    \]
    for some (not necessarily unique) elements $y_i \in \gamma_i(G)$ for $k \leq i \leq n$. The second claim follows by inspecting each term of this product as above. 
\end{proof}

\end{document}
